\section{Timing: protocol, threads and an earlier measurement}\label{supp:timing}

% generated from the committed records; do not edit
\paragraph{The measurements of \cref{tab:cost}}
The two-dimensional column times the first 32 validation instances of the two-dimensional comparison, and the three-dimensional columns the first 32 validation and 8 fine-mesh instances of the three-dimensional comparison, with the label-free transformers of seed 0. The table gives medians over the instances of the wall time per instance, all load cases together and without the assembly, which the surrogate does not need; the iterations of conjugate gradients (CG), from the zero field and to the prediction's accuracy, and CG's displacement error are medians over the load cases. The surrogate ran on one NVIDIA GeForce RTX 5090 with TF32, preparing the inputs and copying the prediction back included; its first call was timed after an untimed warm-up on 4 validation instances, and its repeated call is the median of 2 further passes. The solvers ran on the CPU of the same machine (Xeon Platinum 8470Q, 25 CPUs, the thread environment set to 25): SuperLU, which produced the labels, and unpreconditioned CG to a relative residual of $10^{-10}$, from the zero field or from the prediction. CG to the prediction's accuracy stops, on each load case, at the first iterate whose relative energy error is at most the prediction's; it reads the reference solution, so that a practical run would take at least as long. The fine-mesh direct solve used up its budget of 900 s after seven of the eight instances; its median, and the ratios that use it, are over those seven. The Ritz projection ran from the lowest to the highest degree ($p = 6$ to 9 in three dimensions, 6 and 8 in two) with the machine's 25 BLAS threads; its time is that after the assembly of $K$, which the table lists apart (in two dimensions the assembler also builds the four load vectors), for the free block, the basis, $\Phi^{\top}K\Phi$ and $\Phi^{\top}F$, the small solve and the reconstruction. On the fine mesh, three of the 32 load cases had a relative energy gap above 1 (4.76, 4.77 and 22.0) and three of the eight instances a mean above 1; CG matches such a load case with no iteration.


\paragraph{Threads}

The solvers of \cref{tab:cost} ran with the machine's \numCostThreads{} threads. The effect
of the number of threads was measured on another host of the same type, with
\numThreadsHost{} threads, on the first \numThreadsValN{} validation and \numThreadsFineN{}
fine-mesh instances of the three-dimensional comparison: with one thread, the direct solve
took the same time, to within \numThreadsDirectDev{}, and conjugate gradients
\numThreadsCgGainMin{} to \numThreadsCgGainMax{} less; a repeat with all threads agreed with
the first measurement to within \numThreadsDrift{}. In two dimensions (\numThreadsTwoN{}
instances, without a repeat), one thread was \numThreadsTwoDirect{} faster for the direct
solve and \numThreadsTwoCg{} for conjugate gradients.

\paragraph{An earlier measurement}
The timings of \cref{tab:cost} were measured on 28 October 2026, on the machine
that computed the Ritz projection, so that every row of the table comes from one machine.
The same instances and networks had been timed on 8 October 2026 on another machine
(\numRerunCpu{}, \numRerunThreads{} threads). The surrogate's accuracies agreed between the
two measurements to within a relative \numRerunAccDev{}, and the iteration counts of
conjugate gradients differed by at most \numRerunIters{} per load case. The times of the
earlier measurement are committed with the records.
